% =============================================================================
% subselect — methodology log
%
% This file records finalized methodology decisions for the subselect framework.
% It is the source-of-truth that the EGU 2026 poster and the follow-up paper
% will both quote. See ../docs/documentation_workflow.md for what goes here
% and what doesn't.
%
% Compile: pdflatex methods.tex
% =============================================================================

\documentclass[11pt,a4paper]{article}

\usepackage[utf8]{inputenc}
\usepackage[T1]{fontenc}
\usepackage{lmodern}
\usepackage{microtype}
\usepackage{geometry}
\geometry{margin=2.2cm}
\usepackage{amsmath, amssymb, amsthm}
\usepackage{graphicx}
\usepackage{hyperref}
\usepackage[round]{natbib}
\usepackage{enumitem}
\setlist{itemsep=2pt, topsep=4pt}

\title{\texttt{subselect} --- Methodology Log}
\author{Athanasios Tsilimigkras \and Aristeidis Koutroulis \\
        Technical University of Crete, Chania, Greece}
\date{Last updated: \today}

\begin{document}
\maketitle

\begin{abstract}
\noindent
This document records the finalized methodology decisions of the \texttt{subselect}
framework for country-scale CMIP6 sub-ensemble selection. Each decision lists the
question, the alternatives considered, the chosen option, the rationale, the
defining equation or pseudocode, the implementation pointer, and the date.
The document is appended to as decisions are made, never edited retroactively.
It feeds the EGU~2026 poster \citep[abstract EGU26-14502]{subselect:abstract}
and the follow-up framework paper.
\end{abstract}

\tableofcontents

% =============================================================================
\section{Historical performance}\label{sec:performance}
% =============================================================================

The historical-performance dimension is reused unchanged from the published Greece
paper. See \texttt{../docs/historical\_performance.md} for the complete specification
of the variables, reference dataset, evaluation window, and aggregation choices.

% --- entries appended below as performance methodology evolves -----------------

% =============================================================================
\section{Future-response spread}\label{sec:spread}
% =============================================================================

The future-spread dimension is partially inherited from the published paper
(change variables, scenario default) and partially extended for the framework
(formal sub-ensemble coverage metric). See \texttt{../docs/future\_spread.md}
for the design space.

% --- entries appended below as spread methodology is finalized ----------------

% =============================================================================
\section{Model independence}\label{sec:independence}
% =============================================================================

The independence dimension is new in the framework; it was not part of the
published Greece paper. Two methods are scoped: feature-space clustering on a
regional climate-response embedding, and pairwise-RMSE genealogy clustering.
See \texttt{../docs/model\_similarity.md} for the design.

% --- entries appended below as independence methodology is finalized ---------

% =============================================================================
\section{Multi-objective optimization}\label{sec:optimization}
% =============================================================================

The selection engine combines the three dimensions into a scalar cost (default)
and a Pareto front (diagnostic). See \texttt{../docs/cost\_function.md} for the
design space and \texttt{subselect.optimize} for the implementation.

% --- entries appended below as optimization methodology is finalized ---------

% =============================================================================
\section{Country cropping}\label{sec:cropping}
% =============================================================================

The framework supports both bounding-box and shapefile-based cropping of the
input fields. See \texttt{../docs/refactor.md} for the rationale and the
\texttt{subselect.geom} module for the implementation.

% --- entries appended below as cropping policy is finalized ------------------

\subsection{Framework default crop rule}\label{sec:cropping-default}

\paragraph{Question.} Which cropping rule should the framework apply by default
when extracting a CMIP6 model's data over a user-specified country?

\paragraph{Alternatives considered.}
\begin{itemize}
  \item \textbf{\texttt{bbox}}: rectangular slice from
        \texttt{country\_codes.json}, no shapefile dependency. Paper-era
        setting; over-includes ocean and neighbouring land, especially for
        archipelagic countries like Greece.
  \item \textbf{\texttt{shapefile\_strict}}: binary mask via
        \texttt{rasterio.features.geometry\_mask} with
        \texttt{all\_touched=False} --- a pixel is included only if its
        centroid lies inside the country polygon. Risk: leaves zero pixels
        for small countries on coarse grids (e.g.\ Cyprus on CanESM5
        $\sim$2.8\,\textdegree).
  \item \textbf{\texttt{shapefile\_lenient}}: binary mask via
        \texttt{rasterio.features.geometry\_mask} with
        \texttt{all\_touched=True} --- any pixel touched by the country
        polygon is included. NaN outside the country.
  \item \textbf{\texttt{shapefile\_fractional}}: per-pixel area-fraction
        inside the country polygon, in $[0, 1]$, via per-cell
        \texttt{shapely} intersection. Returns the data un-masked plus a
        weight array. Boundary-precision option; opt-in.
\end{itemize}

\paragraph{Choice.} \texttt{shapefile\_lenient}.

\paragraph{Rationale.} Three reasons. First, \emph{operational simplicity}:
the cropped result is a rectangular \texttt{DataArray} with NaN outside the
country, which xarray and matplotlib both handle natively; the existing
$\cos(\varphi)$ area-weighted spatial-mean code works unchanged because the
\texttt{xr.DataArray.weighted(\dots).mean(\dots)} primitive already ignores
NaN entries. Second, \emph{prototype consistency}: the legacy
\texttt{climpact/shp\_extraction.ipynb} prototype already uses this rule, and
adopting it as the default avoids re-litigating an empirical choice that was
satisfactory in earlier exploratory work. Third, \emph{empirical validation}:
the M5 review (visual comparison of all four methods over Greece, Cyprus, and
Switzerland on CanESM5 and MPI-ESM1-2-HR; see
\texttt{notebooks/exploratory/m5\_cropping\_review/}) confirmed that
\texttt{shapefile\_lenient} captures the country footprint sensibly across
both archipelagic (Greece), small-island (Cyprus), and mountainous-landlocked
(Switzerland) geometries on coarse and fine native grids. The paper's
regression test continues to pin \texttt{bbox} for backward compatibility;
\texttt{shapefile\_fractional} remains available for boundary-precision
use cases that justify the larger interface change to weight-aware metrics.

\paragraph{Definition.} For each pixel $p$ on the model's native grid with
spatial extent $A_p$:
%
\begin{equation}
  m(p) =
  \begin{cases}
    1 & \text{if } A_p \cap \Omega \neq \emptyset, \\
    0 & \text{otherwise},
  \end{cases}
\end{equation}
%
where $\Omega$ is the country polygon from GADM 4.1 (layer \texttt{ADM\_0}).
The mask is computed against the full input grid; both the data and the mask
are then cropped to a mask-aware bounding box (the smallest lat/lon box
enclosing all $m(p)=1$ pixels, padded by one full cell on each side). The
cropped data is masked with NaN where $m(p) = 0$:
%
\begin{equation}
  d_{\text{crop}}(p) =
  \begin{cases}
    d(p) & \text{if } m(p) = 1, \\
    \text{NaN} & \text{if } m(p) = 0.
  \end{cases}
\end{equation}

\paragraph{Implementation.} \texttt{subselect.geom.crop} with
\texttt{method="shapefile\_lenient"} (the default kwarg value); pinned by
\texttt{tests/test\_geom.py::test\_lenient\_includes\_cell\_when\_polygon\_lands\_inside\_a\_single\_cell}
and the overshoot-half-cell regression case alongside it.

\paragraph{Decided.} 2026-04-30, package version 0.1.0.dev0.

% =============================================================================
\section{Caching and reproducibility}\label{sec:caching}
% =============================================================================

The framework caches per-country results in three layers: parquet for tabular
metric tables, zarr for fields and matrices, and a SQLite catalog as an index.
See \texttt{../docs/refactor.md} for the design.

% --- entries appended below as caching/reproducibility decisions are made ----

\subsection{Cache stack: parquet + zarr + SQLite}\label{sec:caching-stack}

\paragraph{Question.} Which storage stack should back the per-country and
ensemble-wide derived artefacts (HPS / BVS tables, change signals, climatology
fields, GWL crossing years, distance matrices)?

\paragraph{Alternatives considered.}
\begin{itemize}
  \item \textbf{Parquet (tabular) + Zarr (arrays) + SQLite (index).}
        Three layers, each with one responsibility. Parquet is language-portable,
        schema-explicit, snappy-compressed; Zarr v2 directory stores serve
        trivially as static files; SQLite is in the Python stdlib.
  \item \textbf{DuckDB on Parquet.} Same parquet files, DuckDB replaces SQLite
        as the catalog \emph{and} runs analytical queries directly against
        the parquet files. Extra dependency.
  \item \textbf{Intake catalog} (\texttt{intake-xarray}). YAML-based,
        declarative, climate-community standard for indexing xarray datasets.
  \item \textbf{HDF5 hierarchies} (\texttt{h5py} or \texttt{netCDF4}).
        Single-file, hierarchical container.
\end{itemize}

\paragraph{Choice.} Parquet (tabular) + Zarr (arrays) + SQLite (index).

\paragraph{Rationale.} Four reasons. First, \emph{workload fit}: tabular
outputs are KB to MB-scale per artefact (the canonical Greece directory is
$\sim$138\,MB across 28 spreadsheets), and the hot path is per-country
DataFrame loads in the Phase~2 cost-function inner loop. DuckDB's value is
analytical SQL across many parquet files; that workload does not exist
here, so SQLite as a pure path-catalog is the right weight class. Second,
\emph{consumer set}: the cache is read only by \texttt{subselect} itself and
the future static web app; \texttt{intake-xarray}'s value is sharing a
declarative catalog with third-party climate tooling, which is not a
requirement. Third, \emph{web-app delivery}: zarr~v2 directory stores serve
trivially over HTTP as static files (the Phase~4 web-app delivery model);
HDF5 single-files require a server runtime to slice. Fourth,
\emph{reproducibility hooks}: cache invalidation by code version is a
single SQLite column; harder to retrofit into intake YAML or HDF5 attributes.
The legacy code's pickle artefacts (in \texttt{legacy/cmip6-greece/}) are
language-locked and version-fragile; parquet replaces them and removes
that dependency.

\paragraph{Definition.} The catalog schema enforces the natural key
$(\texttt{country}, \texttt{scenario}, \texttt{season}, \texttt{kind},
\texttt{crop\_method}, \texttt{config\_hash})$ as \texttt{UNIQUE}. Sentinel
values \texttt{\_global} (country) and \texttt{n/a} (scenario / season /
crop\_method) are used for global and time-series artefacts so the
\texttt{NOT NULL} constraint stays enforceable. Three parquet path
conventions:
%
\begin{align*}
  \text{per-country} \quad &\text{cache/parquet/}\langle c \rangle / \langle s \rangle / \langle p \rangle / \langle k \rangle\_\_\langle m \rangle \texttt{.parquet} \\
  \text{time-series} \quad &\text{cache/parquet/}\langle c \rangle / \texttt{timeseries} / \langle v \rangle\_\_\langle m \rangle \texttt{.parquet} \\
  \text{global}       \quad &\text{cache/parquet/\_global/} \langle k \rangle \texttt{.parquet}
\end{align*}
%
where $c$ = country, $s$ = scenario, $p$ = season, $k$ = kind, $m$ =
crop method, $v$ = variable. Zarr stores live at
\texttt{cache/zarr/}$\langle c \rangle$/$\langle s \rangle$/$\langle a
\rangle$\_\_$\langle m \rangle$\texttt{.zarr}.

\paragraph{Implementation.} \texttt{subselect.cache.Catalog} for the SQLite
index plus \texttt{write\_parquet}, \texttt{read\_parquet},
\texttt{write\_parquet\_timeseries}, \texttt{read\_parquet\_timeseries},
\texttt{write\_parquet\_global}, \texttt{read\_parquet\_global},
\texttt{write\_zarr}, \texttt{read\_zarr}, and \texttt{export\_to\_xlsx}.
\texttt{zarr<3} pinned in \texttt{environment.yml} and
\texttt{pyproject.toml} for v2 directory-store stability.

\paragraph{Decided.} 2026-05-01, package version 0.1.0.dev0.

% =============================================================================
\bibliographystyle{plainnat}
\bibliography{references}
% =============================================================================

\end{document}
